\chapter{1928:Adolf Windaus}

\label{chap:chemistry1928}

The Nobel Prize in Chemistry for the year 1928 was awarded to Adolf Windaus.

\textbf{Motivation:}

for the services rendered through his research into the constitution of the sterols and their connection with the vitamins

\section{Adolf Windaus}

\subsection*{Early Life and Education}

Adolf Otto Reinhold Windaus was born on 1876-12-25 in Berlin, Germany.
He was a German organic chemist awarded the Nobel Prize in Chemistry in 1928 for his research into the constitution of the sterols and their connection with the vitamins.

\subsection*{Career and Major Research}

Windaus began as a medical student but turned to chemistry and received his doctorate in 1899 at the University of Freiburg under Heinrich Kiliani, working on the Digitalis compounds.
He worked as an assistant to Emil Fischer in Berlin before returning to Freiburg, where he habilitated in 1903 with a thesis on cholesterol.
He held professorships in Freiburg and at the University of Innsbruck from 1913 to 1915, and in 1915 he succeeded Otto Wallach as professor of chemistry at the University of Göttingen, where he worked until his retirement in 1944.
In 1919 he demonstrated the chemical relationship between cholesterol and the bile acids.
With Alfred Fabian Hess he found that ultraviolet irradiation gives ergosterol antirachitic activity, and the photochemically formed product became known as vitamin D2 (ergocalciferol).
His institute in Göttingen also investigated the antineuritic vitamin, vitamin B1.

\subsection*{Nobel Prize-winning Work}

The work for which Adolf Windaus was awarded the Nobel Prize in Chemistry in 1928 was described by the Nobel Committee as: "for the services rendered through his research into the constitution of the sterols and their connection with the vitamins".

Windaus's research established the constitution of the sterols, beginning with cholesterol, and showed their relationship to the bile acids.
His investigations connected the sterol field to the chemistry of the vitamins, since he proved that irradiation of ergosterol produces the antirachitic vitamin D.

\subsection*{Significance and Impact}

Windaus's work placed the sterols at the center of the chemistry of natural products and made possible the industrial production of vitamin D by the irradiation of ergosterol.
His structural proposals for cholesterol, presented in his Nobel lecture, were later refined with the aid of X-ray crystallography.

\subsection*{Later Career and Honors}

Windaus retired in 1944 but continued to live in Göttingen.
He received the Pour le Mérite for Sciences and Arts, the Grand Cross of the Order of Merit of the Federal Republic of Germany, and honorary doctorates from several universities.
He died on 1959-06-09 in Göttingen, Germany.

\subsection*{Legacy}

The legacy of Adolf Windaus endures through the lasting impact of his scientific contributions.
The chemistry of the sterols and the vitamins that he founded continues to be built upon by researchers across the chemical sciences.

\subsection*{Modern Developments}

Windaus's structural work on cholesterol and steroids, including the role of ultraviolet light in forming vitamin D, led to the modern understanding and prevention of rickets. Vitamin D research now connects to bone health, immunity, and many areas of medicine. Windaus's steroid chemistry also supported the later elucidation of bile acids and sex hormones.

\subsection*{Selected Publications}

\begin{itemize}

\item Windaus, A. (1919). "Die Konstitution des Cholesterins." Nachrichten der Gesellschaft der Wissenschaften zu Göttingen.

\item Windaus, A. (1931). "Antirachitisches Vitamin aus bestrahltem Ergosterin." Hoppe-Seylers Zeitschrift für physiologische Chemie, 203, 70--75.

\end{itemize}

\clearpage

\section*{Chapter Summary}

The 1928 prize recognized Adolf Windaus for his research into the constitution of the sterols and their connection with vitamins. He helped establish the steroid skeleton and showed that irradiation of ergosterol produces vitamin D.
